\printedpage{59}
% head: A. Grothendieck : Dérivateurs, Ch. XVII, Édition des Universités de Montpellier II et Paris VII
% foot: Version 9 février 2019

Now the relation in question $h_0\tilde{i} = \tilde{h}_0$ is equivalent to the set of the two relations
\[
\begin{aligned}
(h_0\tilde{i})r &= h_0(\tilde{i}r)\\
&= h_0i_0\\
&\overset{?}{=} \alpha_0\\
&= \tilde{h}_0r,
\end{aligned}
\]
equivalent to $h_0i_0 = \alpha_0$,
\[
\begin{aligned}
(h_0\tilde{i})s &= h_0(\tilde{i}s)\\
&= h_0\lambda_B\\
&\overset{?}{=} \lambda_Xh'_{I'}\\
&= \tilde{h}_0s,
\end{aligned}
\]
equivalent to $h_0\lambda_B = \lambda_Xh'_{I'}$, that's OK.

This being so, in the square (6.3.14), the hypotheses $i \in \Phi(I)$, $p \in \Psi\langle I\rangle$ imply $\tilde{i}_0 \in \Phi$, $p_o \in \Psi$. Thus the desired arrow $h_0$ exists, q.e.d. % [sic] p_o -> p_0, index printed as letter o, elsewhere index 0 (h_0, i_0)

This completes the proof of corollary 6.3, and thereby also, that of the last part, d), of 6.1. We are happy\,!

\msp{page 117}

\textbf{Proposition 6.4.}\footnote[34]{Subject to the existence of the relevant $\varinjlim$ and $\varprojlim$ in $\mathcal{M}$. We may assume for simplicity $\mathcal{M}$ stable under $\varinjlim$ and $\varprojlim$ .} % mark printed as superscript (34)
\textit{Let $\mathcal{M}$ be a category, and $\Phi \longleftrightarrow \Psi$ a Q-closed couple in Q-duality ($\Phi = \Psi^*$, $\Psi = \Phi_*$). Let $u : I \longrightarrow I'$ be an arrow in} Cat\textit{, whence a sequence of adjoint functors}
\[
\begin{tikzcd}[column sep=large]
\mathcal{M}(I) \arrow[r, shift left=3, "u_!"] \arrow[r, shift right=3, "u_*"'] & \mathcal{M}(I') \arrow[l, "u^*"']
\end{tikzcd}\ .
\]
\textit{We define $\Phi(I)$, $\Psi\langle I\rangle$, $\Phi(I')$, $\Psi\langle I'\rangle$ as in} 6.2.5. \textit{We then have:}

\cercle{a}
\[
\left\{
\begin{aligned}
u_!(\Phi(I)) &\ \subset\ \Phi(I')\ ,\\
u^*(\Psi\langle I'\rangle) &\ \subset\ \Psi\langle I\rangle\ .
\end{aligned}
\right.
\tag*{(6.4.1)}
\]

\cercle{b} \textit{We have the equivalence}
\[
u^*(\Phi(I')) \subset \Phi(I) \qquad\Longleftrightarrow\qquad u_*(\Psi\langle I\rangle) \subset \Psi\langle I'\rangle\ .
\tag*{(6.4.2)}
\]

\cercle{c} \textit{The equivalent relations} (6.4.2) \textit{hold in each of the following three cases:}
\begin{enumerate}[label={}, leftmargin=2.5em]
\item 1°) \textit{$u$ is coétale.}
\item 2°) \textit{$I$ is discrete.}
\item 3°) \textit{For every $i' \in \operatorname{Ob} I'$, the connected components of $I/i'$ admit a final object.}
\end{enumerate}

\newpage
\printedpage{60}
% head: A. Grothendieck : Dérivateurs, Ch. XVII, Édition des Universités de Montpellier II et Paris VII
% foot: Version 9 février 2019

\textbf{Corollary 6.5.} \textit{Let $(W,C,F)$ be a special closed Quillen triple in $\mathcal{M}$, and let us define, for every $I$ in} Cat\textit{, $C(I)$, $TC(I)$ and $F(I)$, $TF(I)$ as in theorem} 6.1. a)\textit{. We then have}

\msp{page 118}

\textit{the equivalences}
\[
\Bigl(u^*(C(I'))\ \subset\ C(I)\Bigr) \quad\Longleftrightarrow\quad \Bigl(u_*(TF)\ \subset\ TF'\Bigr)\ ,
\tag*{(6.5.1)}
\]
\[
\Bigl(u^*(TC(I'))\ \subset\ TC(I)\Bigr) \quad\Longleftrightarrow\quad \Bigl(u_*(F)\ \subset\ F'\Bigr)\ .
\tag*{(6.5.2)}
\]
\textit{The equivalent conditions} (6.5.1) \textit{as well as} (6.5.2) \textit{are satisfied when we are in one of the cases} 1°) 2°) 3°) \textit{of} 6.4. c).

\textsc{Proof} of 6.4. The assertions \cercle{a} , \cercle{b} are $\pm$ tautological, by adjunction, compare XIII, prop. 2, and cor., page 46. To prove c), it suffices to treat case 3°), for 1° $\Longrightarrow$ 3°, 2° $\Longrightarrow$ 3°. But the computation of
\[
u_*(X)(i')\ \simeq\ \varprojlim_{i\in I/i'} X(i)
\]
shows that if the connected components $(I/i')_\alpha$ of $I/i'$ each have a final object $e_\alpha$, then
\[
u_*(X)(i')\ \simeq\ \prod_\alpha X(e_\alpha)\ .
\]
Now $\Psi$, being right closed, is stable under arbitrary products, so $X \in \Psi\langle I\rangle$ implies that the $u_*(X)(i')$ are in $\Psi$, so $u_*$ is in $\Psi\langle I'\rangle$, q.e.d.